\subsection{Earlier approaches to strongly polynomial time}
\label{sec:lp-context}

The first polynomial-time algorithms for general linear programming
measure the size of a rational instance in bits. Khachiyan's ellipsoid
algorithm~\cite{Khach,Khach2} draws on the cutting-plane methods of
Shor and Yudin--Nemirovskii~\cite{Shor,YN}. In its feasibility argument,
the volume of the enclosing ellipsoid falls by a controlled factor at
each cut. Determinant bounds in terms of the input length tell the
algorithm when that volume is too small to contain a feasible region
of the required thickness~\cite{GLS}. Karmarkar's projective
algorithm~\cite{Kar} gave a second route to polynomial bit complexity.
Subsequent interior-point methods follow solutions of penalized problems
as the penalty for approaching the boundary decreases~\cite{NN}.
Faster linear algebra reduces the work per solve~\cite{Ren,Vai,CLS,VDB},
while exact recovery still has to account for the input
precision~\cite{GPV}. Even on the interval $[0,1]$, the objective
$x\mapsto 2^{-L}x$ gives the two vertices values only $2^{-L}$ apart.
Intervals of radius $2^{-k}$ around these values are disjoint only for
$k>L+1$. A single exact sign comparison selects the optimal endpoint.

\paragraph{Removing dependence on the data.}
Tardos first obtained a strongly polynomial algorithm for minimum-cost
circulation~\cite{TarCirc}, then showed that the arithmetic operation
count for a rational LP can depend on the encoding length of the
constraint matrix alone~\cite{Tardos}. This gives a strongly polynomial
bound whenever the entries of the matrix have polynomial encoding length
in the dimensions, including matrices with entries in $\{0,1,-1\}$.
Frank and Tardos~\cite{FT} replace a vector by a short integer vector
whose inner products with bounded integer vectors have the same signs.
When the constraint matrix has short entries, this lets one compress
the objective and the right-hand side while preserving the comparisons
needed for optimization~\cite[\S6.6]{GLS}.

A different source of strongly polynomial bounds is structure in the
constraint matrix. Megiddo solved LP in linear time for fixed
dimension~\cite{MegLin}, and solved feasibility for inequalities
involving at most two variables each~\cite{Meg}. Generalized flow
maximization supplies a related class with at most two nonzeros per
column~\cite{Vegh,OV}. Dadush, Koh, Natura, Olver and V\'egh obtained a
strongly polynomial optimization algorithm for every matrix with at
most two nonzeros in each row or at most two in each
column~\cite{DKNOV}. Their later work treats classes for which the
central path has polynomial straight-line complexity~\cite{DKNOV26}.

\paragraph{Geometry of the constraint matrix.}
For a real matrix $A$, Vavasis and Ye~\cite{VY} gave a path-following
algorithm whose iteration bound depends on the dimensions and the
logarithm of a condition measure of $A$. Their layered least-squares
steps can cross several numerical scales at once. A related condition
measure in standard form $Ax=b$, $x\ge0$ has a concrete description.
A \emph{circuit} of $A$ is a nonzero kernel vector of minimal support.
The \emph{circuit imbalance} is the largest ratio of magnitudes of
two nonzero coordinates in the same circuit~\cite{Rock,CIRC}.
Dadush, Huiberts, Natura and V\'egh~\cite{DHNV} obtained an iteration
bound in terms of the best condition achievable by diagonal scaling of
the columns. Circuit imbalance also controls the precision needed in
an approximate-to-exact solver~\cite{DNV}.

Even the best rescaling can leave large ratios in a fixed dimension.
For $0<u<1$, the matrix
\[
 A=\begin{pmatrix}1&0&1&u\\0&1&u&1\end{pmatrix}
\]
has circuits $(-1,-u,1,0)$ and $(-u,-1,0,1)$. The ratios of the
magnitudes of their first two coordinates are $1/u$ and $u$.
Any positive diagonal column scaling multiplies these two ratios by the
same factor, so at least one ratio or its reciprocal is at least $1/u$.
A separate obstruction comes from the central path itself. Allamigeon,
Benchimol, Gaubert and Joswig constructed instances requiring
exponentially many segments to follow the logarithmic central path
in a prescribed neighborhood~\cite{ABGJ}. Allamigeon, Gaubert and
Vandame extended this obstruction to path-following with
self-concordant barriers~\cite{AGV}.

\paragraph{Edge paths and rescalings.}
Early applications ranged from Kantorovich's production planning~\cite{Kan}
to the computation of optimal strategies in zero-sum matrix games by LP
duality~\cite{Gale}. Dantzig's simplex method exploits the geometry of these
problems by moving between adjacent feasible bases~\cite{Dantzig}.
Klee and Minty exhibited exponentially long runs for a classical
choice of pivot~\cite{KM}. Randomized pivot rules achieve
subexponential bounds~\cite{Kalai,MSW}, and smoothed analysis gives
polynomial expected bounds in the dimensions and the inverse noise
scale~\cite{ST,BH}. Awoniyi has proposed a different pivot-based
approach to strongly polynomial LP~\cite{Awon}.

Even the shortest possible edge walk can be long. Kalai and
Kleitman proved a quasipolynomial diameter bound for polyhedral
graphs~\cite{KK}, and Santos disproved the original Hirsch bound of
the number of facets minus the dimension~\cite{Santos}. The companion
paper~\cite{OAIHirsch} constructs simple rational polytopes whose
graph diameter grows faster than any polynomial in their number of
facets. A circuit walk can leave the edge graph by taking a maximal
feasible step in a circuit direction. Natura constructs polynomial-length
monotone circuit walks between specified vertices~\cite{Natura}.

The algorithm of Theorem~\ref{thm:lp} searches over coordinate
rescalings. In earlier rescaling algorithms, finding that one coordinate
must be small in every normalized feasible point can guide a change
of scale that improves the geometry~\cite{Betke,DV,Chu,DVZ,PS}. Here a
rescaling selects the coordinates on which a primal or dual witness
can be detected. Section~\ref{sec:overview} explains how the algorithm
combines such observations, even when the initial system has arbitrarily
large circuit imbalance.
